\section{Spectral localizers in finite and periodic systems}
\label{sec:localizer}
\subsection{Energy-position tuples and Clifford structure}
A spectral localizer combines operators associated with several
observables into one self-adjoint matrix. For two positions and a
Hamiltonian, Pauli matrices provide a Clifford representation.
In an orthonormal space define
\begin{equation}
 \widetilde L=
 (\widetilde H-EI)\otimes\sigma_z+
 \kappa(\widetilde X-xI)\otimes\sigma_x+
 \kappa(\widetilde Y-yI)\otimes\sigma_y ,
 \label{eq:clifford-localizer}
\end{equation}
up to the explicitly fixed tensor ordering.
The parameter $\kappa>0$ has units of energy/length.
It balances energy and position resolution rather than correcting
the Hamiltonian. The connection between a localizer signature and
an even index pairing holds under appropriate locality, gap, and
finite-volume bounds \cite{LoringSchulzBaldes2020}.
Those hypotheses remain relevant when interpreting a finite Gaussian model.

In auxiliary-index-first ordering the covariant AO matrix is
\begin{equation}
 L=
 \begin{pmatrix}
 H-ES & \kappa[(X-xS)-\ii(Y-yS)]\\
 \kappa[(X-xS)+\ii(Y-yS)] & -(H-ES)
 \end{pmatrix},\qquad
 B=\begin{pmatrix}S&0\\0&S\end{pmatrix}.
 \label{eq:covariant-localizer}
\end{equation}
For SOC, $H$ and $S$ here denote the physical spinor matrices.
Auxiliary doubling is separate from spin, so the localizer dimension
is four times the scalar AO count. All AO directions are retained:
there is no hidden occupied-space projection or automatic removal
of nearly dependent basis functions.

\subsection{Index, metric gap, and stability}
For class A define
\begin{equation}
 \nu_A=\frac12\Sig(L)=\frac12(n_+-n_-),\qquad
 \mu=\min_j|\lambda_j|,\quad Lv_j=\lambda_j Bv_j .
 \label{eq:index-gap}
\end{equation}
Signature is invariant under nonsingular congruence, whereas ordinary
eigenvalues of a covariant AO matrix are not.
Inertia can therefore be evaluated directly from $L$, but the
protection gap must retain $B$. Equivalently, $\mu$ is the smallest
absolute eigenvalue of $B^{-1/2}LB^{-1/2}$; this identity does not
require forming that dense matrix.

At fixed metric, the index cannot change under a perturbation satisfying
\begin{equation}
 \|B^{-1/2}\delta L B^{-1/2}\|<\mu .
 \label{eq:localizer-stability}
\end{equation}
This finite-matrix statement gives local protection an operational
meaning. If the metric changes, the corresponding orthonormalized
problems must be compared, including the effect of $\delta B$.
An unresolved near-zero gap yields an \emph{unresolved} index,
not evidence for trivial topology.

Moving the query and coordinate operators together, shifting the
Hamiltonian and query energy together, or changing AO basis must
leave the generalized gap unchanged. Reversing the oriented plane
reverses the class A sign. Model comparisons establish the sign
convention before applying it to a material.

\subsection{Inertia and gap require different calculations}
A pivoted Hermitian factorization has the form
\begin{equation}
 P^\dagger LP=\mathcal L D\mathcal L^\dagger .
 \label{eq:ldl}
\end{equation}
Here $\mathcal L$ is unit lower triangular and $D$ contains
$1\times1$ and $2\times2$ blocks.
Sylvester's law transfers the inertia from $L$ to $D$.
For a Hermitian block,
\[
 D_b=\begin{pmatrix}a&b\\\overline b&d\end{pmatrix},\qquad
 \lambda_\pm=\frac{a+d}{2}
 \pm\sqrt{\left(\frac{a-d}{2}\right)^2+|b|^2}.
\]
Stable evaluation must retain scaling and both signs.
A zero-diagonal block with $b\ne0$ has one positive and one negative
direction; counting its diagonal entries is wrong.
The dense reference uses LAPACK \texttt{ZHETRF}, with \texttt{DSYTRF}
for exactly real matrices \cite{LapackLDL}. A small imaginary part is
not discarded to select real arithmetic.

Pivot magnitudes are not generalized gaps. For $\sigma>0$,
shifted inertia counts satisfy, away from singular endpoints,
\begin{equation}
 \begin{split}
 n_-(L-\sigma B)-n_-(L)&=\#\{j:0<\lambda_j<\sigma\},\\
 n_-(L)-n_-(L+\sigma B)&=\#\{j:-\sigma<\lambda_j<0\}.
 \end{split}
 \label{eq:gap-inertia}
\end{equation}
These counts bracket the nearest positive or negative generalized
eigenvalue. Factorization failure or an ambiguous endpoint must not
be treated as a successful certificate. The reported bounds are
finite-precision numerical brackets, not interval-arithmetic guarantees.

\subsection{Pfaffian orientation and physical time reversal}
For a real skew-symmetric matrix $A$ of even order,
$\det A=\Pf(A)^2$: the determinant loses the Pfaffian sign.
Moreover, $\Pf(P^TAP)=\det(P)\Pf(A)$.
An odd simultaneous row/column permutation reverses that sign
without changing the determinant.
Ordinary inertia cannot replace a skew factorization with explicit
permutation bookkeeping.

For spinful time reversal with $\Theta^2=-1$, the localizer has
a real skew representation \cite{DollSchulzBaldes2021,Wong2026}.
In auxiliary/spin/AO ordering, let $C=J_{\mathrm{aux}}\otimes J$
and $Q=(I-\ii C)/\sqrt2$. The symmetry relations imply that
\begin{equation}
 A=\ii Q^\dagger LQ
 \label{eq:aii-skew}
\end{equation}
is real and skew-symmetric. Both properties are checked numerically,
with residuals small relative to the independently resolved metric
gap. An arbitrary complex AO basis change also transforms time
reversal; keeping $J$ fixed would generally be inconsistent.

Let $A_0$ be the corresponding transform of
$L_0=\operatorname{diag}(S,-S)$. The reference-calibrated index is
\begin{equation}
 (-1)^{\nu_{\mathrm{AII}}}
 =\operatorname{sgn}\Pf(A)\,
  \operatorname{sgn}\Pf(A_0).
 \label{eq:pfaffian-index}
\end{equation}
Both factors use the same ordering and orientation.
The dense algorithm accumulates signs rather than an overflowing
Pfaffian product. This is a time-reversal local index,
not an inversion-parity test.

The native sparse skew factorization implements a separate numerical
task from symmetric inertia. It retains a distributed Schur complement
and oriented $2\times2$ pivots, accumulating the pivot signs and the
parity of their permutations. The transformation in
Eq.~\eqref{eq:aii-skew} is assembled by distributed addition of
coordinate entries rather than by gathering a dense matrix.
Independent right-hand sides test the retained factors against the
assembled matrix. The discarded symmetric and imaginary parts are
bounded relative to the independently resolved metric gap before
an index is accepted. This calculation requires neither MUMPS nor
Tacho. The optional Tacho comparison implements the same separate
skew-Pfaffian task \cite{Wong2026}, but its factorization resides on
one rank and is not an MPI-distributed Pfaffian solver.
The motivating photonic and aperiodic examples
\cite{Dixon2023,Wong2026} provide methodological context, not
first-principles material validation for the CP2K path.

\subsection{Finite three-dimensional class AII}
For three Cartesian positions, the finite AII construction uses
a real determinant rather than the two-dimensional Pfaffian
\cite{Loring2015}. In auxiliary-index-first ordering, define
\begin{equation}
 \begin{split}
 A_3&=I_2\otimes(H-ES)
       -\ii\kappa\sum_{j=1}^{3}\sigma_j\otimes(X_j-x_jS),\\
 B_3&=I_2\otimes S,\qquad
 L_3=\begin{pmatrix}0&A_3\\A_3^\dagger&0\end{pmatrix},\qquad
 \mathcal B_3=\operatorname{diag}(B_3,B_3).
 \end{split}
 \label{eq:finite-three-dimensional-localizer}
\end{equation}
Here $X_j$ are the covariant Cartesian moments, $x_j$ are the
query coordinates, and the Pauli matrices $\sigma_j$ act only
in the auxiliary space. The physical-spin Hamiltonian, metric
and all three positions must obey the fixed fermionic
time-reversal representation. With $C=J_{\mathrm{aux}}\otimes J$
and $Q=(I-\ii C)/\sqrt2$, the chiral block
$R_3=Q^\dagger A_3Q$ is real but need not be symmetric.
For a resolved gap, the local index is
\begin{equation}
 \nu_3=\frac{1-\operatorname{sgn}\det R_3}{2},\qquad
 \mu_3=\min_\ell|\lambda_\ell|,\qquad
 L_3v_\ell=\lambda_\ell\mathcal B_3v_\ell .
 \label{eq:finite-three-dimensional-index}
\end{equation}
Equivalently, $\mu_3$ is the smallest singular value of
$B_3^{-1/2}A_3B_3^{-1/2}$. Nonsingular AO congruences
multiply the determinant by a positive factor and preserve
the metric gap. A basis change must also preserve or
consistently transform the time-reversal representation.

The dense implementation accumulates triangular-pivot and
permutation signs without forming a determinant product.
The factor reconstruction residual and a rounding allowance
are compared with the independently calculated metric gap.
If excessive pivot growth prevents acceptance of the triangular
factorization, a Householder factorization supplies orthogonal
and upper triangular factors instead. The determinant sign
includes every nonidentity reflection. The discarded imaginary
part of $R_3$ is included in the perturbation budget, and an
unresolved calculation returns no topological label.

The Hermitian pencil has order eight times the scalar AO count.
Its current solver is a dense reference, including
GPW/GAPW and post-SCF GTH SOC. The periodic extension is given
below, while a sparse three-dimensional solver is not provided.
This strong local index is
not a collection of weak or two-dimensional slice indices.
Agreement with a bulk invariant requires separate volume,
basis and scale convergence. The Supporting Information
includes an explicit finite-box counterexample to automatic
identification with the bulk parity.

\subsection{Periodic positions and the torus construction}
A Cartesian coordinate truncated to a primitive cell is discontinuous
at the periodic seam. Using it in the open-boundary localizer is not
a periodic formulation.
Instead, analytic periodic functions are used \cite{Doll2025}.
For analysis-torus reciprocal vectors $\bG_j$ and
$\phi_j=\bG_j\cdot\br_0$, define
\begin{equation}
 \begin{split}
 C_j&=\langle\phi_\mu|\cos(\bG_j\cdot\br-\phi_j)|\phi_\nu\rangle,\\
 T_j&=\langle\phi_\mu|\sin(\bG_j\cdot\br-\phi_j)|\phi_\nu\rangle .
 \end{split}
\end{equation}
In the implemented orientation,
\begin{equation}
 L_{\mathrm{per}}=
 \begin{pmatrix}
 H-ES-\eta(2S-C_1-C_2)&\eta(T_1-\ii T_2)\\
 \eta(T_1+\ii T_2)&-H+ES+\eta(2S-C_1-C_2)
 \end{pmatrix}.
 \label{eq:periodic-localizer}
\end{equation}
The energy scale $\eta$ is distinct from $\kappa$, and the metric
remains $\operatorname{diag}(S,S)$.
An overall sign and auxiliary-space conjugation relative to the
cited periodic construction make its orientation consistent with
the finite localizer used here.

The trigonometric operators are integrated directly, not evaluated
as matrix functions of finite projected Cartesian positions.
Projection need not preserve either exact unitarity of exponentials
or their commutativity. The theorem for a gapped finite-range lattice
Hamiltonian therefore does not automatically become a theorem for
an arbitrary Gaussian projection.
The output is an AO-torus diagnostic requiring independent
basis-, torus-, and scale-converged material checks.

Explicit analysis k-points describe the enlarged torus through its
complete mesh. Periodic exponentials shift between mesh sectors,
so their cross-k links must be retained.
A symmetry-reduced SCF can supply the potential, but scalar SCF
weights cannot replace the full torus operator.
For three periodic directions, the corresponding chiral block is
\begin{equation}
 \begin{split}
 A_{3,\mathrm{per}}={}&I_2\otimes
   \left[H-ES-\eta\left(3S-\sum_{j=1}^3C_j\right)\right]\\
 &-\ii\eta\sum_{j=1}^3\sigma_j\otimes T_j .
 \end{split}
 \label{eq:periodic-three-dimensional-localizer}
\end{equation}
The phase-shifted $C_j,T_j$ are the analytic AO integrals defined
above, and $\eta>0$ has energy units. Replacing $A_3$ by
$A_{3,\mathrm{per}}$ in Eqs.~\eqref{eq:finite-three-dimensional-localizer}
and \eqref{eq:finite-three-dimensional-index} gives the Hermitian
metric pencil, real determinant sign and gap. In this orientation
the chiral block is an energy-scaled adjoint of the odd-dimensional
periodic construction of Ref.~\cite{Doll2025}. Taking the adjoint
preserves the sign of its real determinant.
The native three-dimensional path uses all three torus directions
with fully periodic cell and Poisson boundaries, fermionic time
reversal and the dense AII solver. Its finite-AO validation includes
skew cells, origin periodicity and GPW/GAPW comparisons. It is not
a sparse implementation or a bulk-convergence certificate.

Spectral flattening provides a separate way to examine the large
bandwidth-to-gap ratio of a DFT Hamiltonian. For positive $S$ and
$E$ in a resolved electronic gap, define
\begin{equation}
 \begin{aligned}
 F_E&=\Delta S\,\operatorname{sign}\!\left[S^{-1}(H-ES)\right]\\
 &=\Delta S^{1/2}\operatorname{sign}\!\left[S^{-1/2}(H-ES)S^{-1/2}\right]S^{1/2},
 \qquad \Delta>0 .
 \end{aligned}
 \label{eq:covariant-flattening}
\end{equation}
This covariant Hermitian operator retains the occupied projector
and replaces generalized band energies by $\pm\Delta$.
The dispersion-reduction argument for periodic localizers motivates
this control \cite{Doll2025}. It does not equate a flattened
finite-volume localizer with the original one: the localizer gap
can close during flattening even when the electronic gap stays open.
An exact sign is also generally longer ranged than $H$.
Consequently, the material tests in the Supporting Information
are distinguished from a proof of the finite-range theorem's
hypotheses and from the native unflattened calculation.

The optional \texttt{SPECTRAL\_FLATTENING} path evaluates
Eq.~\eqref{eq:covariant-flattening} for finite systems with
distributed DBCSR metric inversion and matrix-sign iterations.
Before those iterations,
the original positive metric and the distance of $E$ to the
generalized spectrum of $(H,S)$ are checked by dense diagonalization
or shifted MUMPS factorizations. An unresolved electronic gap is
rejected rather than enlarged by flattening. This precondition is
distinct from the gap of the final localizer. The transformed
Hamiltonian is reused across position and scale queries at fixed
$E$; the SCF Hamiltonian and Gaussian coordinate operators remain
unchanged. The option is off by default. The distributed inverse
and exact-sign approximation can increase fill-in, so this path
does not imply linear scaling.

For a translation-invariant periodic AO torus, the same sign can
instead be evaluated in complete primitive Bloch blocks.
The implementation verifies translation covariance of the entire
Hamiltonian and metric before extracting a representative block row.
It distributes full generalized Bloch eigensystems over MPI ranks,
forms $F_E(k)=\Delta S(k)C(k)\operatorname{sign}(\epsilon_k-E)
C(k)^\dagger S(k)$, and Fourier-transforms back to the original
distributed spin/cell/atom order. The metric bounds and electronic
gap refer to the union of all sampled Bloch spectra. No bands are
discarded and no independent localizer is assigned to a single
k-point. This avoids a whole-torus eigensystem or sign iteration
for the flattening step; it does not remove the final localizer
solve or its possible fill-in.

\subsection{Connection to the quadratic pseudospectrum}
For Hermitian operators $A_j$ and anticommuting Hermitian Clifford
matrices $\gamma_j$, squaring $L=\sum_jA_j\otimes\gamma_j$ gives
\begin{equation}
 L^2=\left(\sum_j A_j^2\right)\otimes I+
 \sum_{j<l}[A_j,A_l]\otimes\gamma_j\gamma_l .
 \label{eq:clifford-square}
\end{equation}
The first term measures simultaneous approximate localization;
the second retains noncommutativity.
The quadratic gap and localizer gap are related, but not
interchangeable. In an AO basis each projected product must
include the metric, or be evaluated after orthonormalization.
